% BHU PYQ -- Manually transcribed & error-corrected
% Paper: BCH-226 : Business Statistics
% Exam: B.Com. (Hons.) Semester IV Examination, 2023-24

\documentclass[11pt,a4paper]{article}
\usepackage[a4paper,top=2.5cm,bottom=2.5cm,left=2.8cm,right=2.8cm,headheight=14pt]{geometry}
\usepackage{amsmath,amssymb}
\usepackage[T1]{fontenc}
\usepackage[utf8]{inputenc}
\usepackage{lmodern,microtype,fancyhdr,enumitem,hyperref,lastpage,parskip}

\hypersetup{
    colorlinks=true,
    urlcolor=black,
    linkcolor=black,
    pdftitle={BCH-226: Business Statistics -- BHU B.Com. (Hons.) Sem IV 2023-24}
}

\pagestyle{fancy}
\fancyhf{}
\renewcommand{\headrulewidth}{0.4pt}
\renewcommand{\footrulewidth}{0.4pt}
\fancyhead[L]{\small\textit{Banaras Hindu University}}
\fancyhead[R]{\small\textit{BCH-226: Business Statistics}}
\fancyfoot[C]{\small Page \thepage\ of \pageref{LastPage}}

\newlist{parts}{enumerate}{2}
\setlist[parts,1]{label=(\roman*),leftmargin=*,itemsep=5pt,topsep=3pt}
\setlist[parts,2]{label=(\alph*),leftmargin=*,itemsep=3pt,topsep=2pt}
\newcommand{\pts}[1]{\hfill\textit{[#1]}}

\begin{document}

\begin{center}
  {\large\bfseries BANARAS HINDU UNIVERSITY}\\[2pt]
  {\normalsize B.Com. (Hons.) --- Semester IV Examination, 2023-24}\\[6pt]
  \rule{0.8\linewidth}{0.5pt}\\[6pt]
  {\large\bfseries Paper No. BCH-226: Business Statistics}\\[4pt]
  {\normalsize Time: 3 Hours\hfill Maximum Marks: 70}
\end{center}

\rule{\linewidth}{0.4pt}

\smallskip
\noindent\textit{Instructions: Answer all questions. All questions carry equal marks (14 marks each).}

\bigskip

\noindent\textbf{Question 1.}
\begin{parts}
  \item Prepare chain base index numbers from the following data:
  \begin{center}
  \small
  \begin{tabular}{|l|c|c|c|c|c|}
  \hline
  \textbf{Year} & 2019 & 2020 & 2021 & 2022 & 2023 \\
  \hline
  \textbf{Price (Rs.)} & 80 & 120 & 132 & 264 & 396 \\
  \hline
  \end{tabular}
  \end{center}\pts{7}
  
  \item From the chain base index numbers given below, prepare fixed base index numbers:
  \begin{center}
  \small
  \begin{tabular}{|l|c|c|c|c|c|c|}
  \hline
  \textbf{Year} & 2018 & 2019 & 2020 & 2021 & 2022 & 2023 \\
  \hline
  \textbf{Chain Base Index No.} & 92 & 102 & 104 & 98 & 103 & 101 \\
  \hline
  \end{tabular}
  \end{center}\pts{7}
\end{parts}

\centerline{\textbf{OR}}

Explain the points to be considered in the construction of Index Numbers.\pts{14}

\medskip\rule{\linewidth}{0.2pt}

\noindent\textbf{Question 2.}
Calculate index numbers for 2020 and 2023 from the following data relating to cost of living of the working class in a city:

\begin{center}
\small
\begin{tabular}{|l|c|c|c|}
\hline
\textbf{Items} & \textbf{Weight} & \multicolumn{2}{c|}{\textbf{Price Index of items}} \\
\cline{3-4}
& & \textbf{2020} & \textbf{2023} \\
\hline
Food & 48 & 110 & 130 \\
Clothing & 8 & 120 & 125 \\
Fuel and Light & 7 & 110 & 120 \\
Room Rent & 13 & 100 & 100 \\
Miscellaneous & 14 & 115 & 135 \\
\hline
\end{tabular}
\end{center}

There is a wage increase of 8\% from the year 2020 to 2023. Is it adequate?\pts{14}

\centerline{\textbf{OR}}

Why Fisher's Index Number is Ideal? Explain with the help of imaginary figures.\pts{14}

\medskip\rule{\linewidth}{0.2pt}

\noindent\textbf{Question 3.}
Calculate trend values by the method of least squares from the data given below:

\begin{center}
\small
\begin{tabular}{|l|c|c|c|c|c|c|c|}
\hline
\textbf{Year} & 2017 & 2018 & 2019 & 2020 & 2021 & 2022 & 2023 \\
\hline
\textbf{Value} & 75 & 67 & 68 & 65 & 50 & 54 & 41 \\
\hline
\end{tabular}
\end{center}\pts{14}

\centerline{\textbf{OR}}

Explain moving average method for measurement of trend with the help of imaginary figures.\pts{14}

\medskip\rule{\linewidth}{0.2pt}

\pagebreak

\noindent\textbf{Question 4.}
Calculate seasonal index from the following data using the average method:

\begin{center}
\small
\begin{tabular}{|l|c|c|c|c|}
\hline
\textbf{Year} & \multicolumn{4}{c|}{\textbf{Quarter}} \\
\cline{2-5}
& \textbf{I} & \textbf{II} & \textbf{III} & \textbf{IV} \\
\hline
2019 & 72 & 68 & 80 & 70 \\
2020 & 76 & 70 & 82 & 74 \\
2021 & 74 & 66 & 84 & 80 \\
2022 & 76 & 74 & 84 & 78 \\
2023 & 78 & 74 & 86 & 82 \\
\hline
\end{tabular}
\end{center}\pts{14}

\centerline{\textbf{OR}}

What are various methods of measuring seasonal variations in time series? Discuss any one method with the help of imaginary figures.\pts{14}

\medskip\rule{\linewidth}{0.2pt}

\noindent\textbf{Question 5.}
\begin{parts}
  \item A total of 10 samples of 8 each was drawn from a universe. The average range ($\bar{R}$) and average standard deviation ($\bar{\sigma}$) of these samples were 0.0811 and 0.0263 respectively. Estimate the standard deviation of the universe from which these samples were drawn when $d_2 = 2.847$ and $c_2 = 0.9027$.\pts{7}
  \item If $\bar{X} = 0.8768$, $\bar{R} = 0.0026$ and given that $A_2 = 0.58$, $D_3 = 0$ and $D_4 = 2.115$. Compute UCL and LCL for $\bar{X}$ and $R$ chart.\pts{7}
\end{parts}

\centerline{\textbf{OR}}

What do you understand by `Statistical Quality Control'? Explain its need and importance in an industry.\pts{14}

\vfill
\rule{\linewidth}{0.4pt}
\begin{center}\small
\href{https://bhu-exam.vercel.app/}{Click here to visit our website}\\[4pt]
\href{https://me-aryan.vercel.app/}{Click here to visit our contributor}\\[4pt]
\href{https://quashan-me.vercel.app/}{Click here to visit our contributor}
\end{center}

\end{document}
